% GENERATED FILE. Edit canonical lessons, then run npm run generate:books.
% canonical-source-hash: fnv1a64:f3ee5188543e11c7
% canonical-lessons: TE-C504-aradhincu, TE-C504-abhimanincu, TE-C504-abhimani, TE-C504-meccuko, TE-C504-prasamsa

\chapter{Admiration and Praise}
\label{ch:te-admiration-and-praise}
\hlchaptermodality{504}

\begin{chapteropening}
\textbf{By the end of this chapter:} \emph{I can say that I admire someone and praise their work.}

Complete the last lesson of chapter 504: I can say that I admire someone and praise their work.
\end{chapteropening}

\section[\texorpdfstring{\te{ఆరాధించు} (ārādhiñcu)}{ఆరాధించు (ārādhiñcu)}]{\te{ఆరాధించు} (ārādhiñcu) — to adore; to worship}

\label{lesson:TE-C504-aradhincu}



\begin{quote}
Before the new one: say the Telugu for to fondle, then the Telugu for fondness.
\end{quote}

\subsection*{You'll want to know: \te{ఆరాధించు}}
\textbf{\te{ఆరాధించు}} — \emph{ārādhiñcu} — \textquotedblleft{}to adore; to worship\textquotedblright{}.

\begin{quote}
\te{ఆ} \te{అమ్మాయి} \te{తన} \te{అక్కను} \te{ఆరాధిస్తుంది}.
\end{quote}

\begin{quote}
\emph{ā ammāyi tana akkanu ārādhistundi.}
\end{quote}

\begin{quote}
That girl adores her elder sister.
\end{quote}

Past (I …): \te{ఆరాధించాను}, \emph{ārādhiñcānu}.

The noun is \te{ఆరాధన}, adoration.

\begin{cousinweb}
\textbf{\te{ఆరాధించు}} (\emph{ārādhiñcu}) is built on Sanskrit \emph{ārādhana}, \textquotedblleft{}worship\textquotedblright{} (CDIAL 1316), the noun of \emph{ārādhayati}, \textquotedblleft{}to win over\textquotedblright{}, with the Telugu verb ending -\te{ఇంచు}.
\end{cousinweb}

\begin{grammarlens}[title={What you've built}]
One more word for describing what you lived and felt.
\end{grammarlens}

\subsection*{Your turn}
\begin{itemize}
\raggedright
  \item \emph{Say it:} \emph{ārādhiñcu}
  \item \emph{Say it:} \emph{ārādhiñcu}, once more
  \item \emph{Say it:} \emph{ā ammāyi tana akkanu ārādhistundi.} — That girl adores her elder sister.
  \item \emph{Recall:} say the Telugu for to fondle, then the Telugu for fondness, then say \emph{ārādhiñcu} again
\end{itemize}

\begin{tcolorbox}[breakable,colback=teal!4,colframe=teal!35!black,title={Before you move on}]
What does \te{ఆరాధించు} mean? (\textquotedblleft{}To adore; to worship\textquotedblright{}.) Say it once more.
\end{tcolorbox}
\section[\texorpdfstring{\te{అభిమానించు} (abhimāniñcu)}{అభిమానించు (abhimāniñcu)}]{\te{అభిమానించు} (abhimāniñcu) — to admire; to be fond of}

\label{lesson:TE-C504-abhimanincu}



\begin{quote}
Before the new one: say the Telugu for fondness, then the Telugu for to adore.
\end{quote}

\subsection*{You'll want to know: \te{అభిమానించు}}
\textbf{\te{అభిమానించు}} — \emph{abhimāniñcu} — \textquotedblleft{}to admire; to be fond of\textquotedblright{}.

\begin{quote}
\te{ఆ} \te{డాక్టరును} \te{అందరూ} \te{అభిమానిస్తారు}.
\end{quote}

\begin{quote}
\emph{ā ḍākṭarunu andarū abhimānistāru.}
\end{quote}

\begin{quote}
Everyone admires that doctor.
\end{quote}

Past (I …): \te{అభిమానించాను}, \emph{abhimāniñcānu}.

The verb of \te{అభిమానం}, regard.

\begin{cousinweb}
\textbf{\te{అభిమానించు}} (\emph{abhimāniñcu}) is built on \te{అభిమానం}, from Sanskrit \emph{abhimāna} (CDIAL 524), with the Telugu verb ending -\te{ఇంచు}.
\end{cousinweb}

\begin{grammarlens}[title={What you've built}]
One more word for describing what you lived and felt.
\end{grammarlens}

\subsection*{Your turn}
\begin{itemize}
\raggedright
  \item \emph{Say it:} \emph{abhimāniñcu}
  \item \emph{Say it:} \emph{abhimāniñcu}, once more
  \item \emph{Say it:} \emph{ā ḍākṭarunu andarū abhimānistāru.} — Everyone admires that doctor.
  \item \emph{Recall:} say the Telugu for fondness, then the Telugu for to adore, then say \emph{abhimāniñcu} again
\end{itemize}

\begin{tcolorbox}[breakable,colback=teal!4,colframe=teal!35!black,title={Before you move on}]
What does \te{అభిమానించు} mean? (\textquotedblleft{}To admire; to be fond of\textquotedblright{}.) Say it once more.
\end{tcolorbox}
\section[\texorpdfstring{\te{అభిమాని} (abhimāni)}{అభిమాని (abhimāni)}]{\te{అభిమాని} (abhimāni) — an admirer; a fan}

\label{lesson:TE-C504-abhimani}



\begin{quote}
Before the new one: say the Telugu for to adore, then the Telugu for to admire.
\end{quote}

\subsection*{You'll want to know: \te{అభిమాని}}
\textbf{\te{అభిమాని}} — \emph{abhimāni} — \textquotedblleft{}an admirer; a fan\textquotedblright{}.

\begin{quote}
\te{నా} \te{తమ్ముడు} \te{ఆ} \te{ఆటగాడికి} \te{పెద్ద} \te{అభిమాని}.
\end{quote}

\begin{quote}
\emph{nā tammuḍu ā āṭagāḍiki pedda abhimāni.}
\end{quote}

\begin{quote}
My younger brother is a big fan of that player.
\end{quote}

Plural: \te{అభిమానులు}, \emph{abhimānulu}.

-\te{కి} \te{అభిమాని}, a fan of.

\begin{cousinweb}
\textbf{\te{అభిమాని}} (\emph{abhimāni}) is borrowed from Sanskrit \emph{abhimānin}, built on \emph{abhimāna} (CDIAL 524): one who holds someone in high regard.
\end{cousinweb}

\begin{grammarlens}[title={What you've built}]
One more word for describing what you lived and felt.
\end{grammarlens}

\subsection*{Your turn}
\begin{itemize}
\raggedright
  \item \emph{Say it:} \emph{abhimāni}
  \item \emph{Say it:} \emph{abhimāni}, once more
  \item \emph{Say it:} \emph{nā tammuḍu ā āṭagāḍiki pedda abhimāni.} — My younger brother is a big fan of that player.
  \item \emph{Recall:} say the Telugu for to adore, then the Telugu for to admire, then say \emph{abhimāni} again
\end{itemize}

\begin{tcolorbox}[breakable,colback=teal!4,colframe=teal!35!black,title={Before you move on}]
What does \te{అభిమాని} mean? (\textquotedblleft{}An admirer; a fan\textquotedblright{}.) Say it once more.
\end{tcolorbox}
\section[\texorpdfstring{\te{మెచ్చుకో} (meccukō)}{మెచ్చుకో (meccukō)}]{\te{మెచ్చుకో} (meccukō) — to appreciate; to praise}

\label{lesson:TE-C504-meccuko}



\begin{quote}
Before the new one: say the Telugu for to admire, then the Telugu for an admirer.
\end{quote}

\subsection*{You'll want to know: \te{మెచ్చుకో}}
\textbf{\te{మెచ్చుకో}} — \emph{meccukō} — \textquotedblleft{}to appreciate; to praise\textquotedblright{}.

\begin{quote}
\te{నాన్న} \te{నా} \te{బొమ్మను} \te{మెచ్చుకున్నారు}.
\end{quote}

\begin{quote}
\emph{nānna nā bommanu meccukunnāru.}
\end{quote}

\begin{quote}
Dad praised my picture.
\end{quote}

Past (I …): \te{మెచ్చుకున్నాను}, \emph{meccukunnānu}.

Warmer than \te{పొగడు}, which can be flattery: \te{మెచ్చుకో} is honest praise.

\begin{cousinweb}
\textbf{\te{మెచ్చుకో}} (\emph{meccukō}) is native Dravidian: DEDR 4722 lists Telugu \emph{meccu}, \textquotedblleft{}to admire, to praise, to approve\textquotedblright{}, in the article of Tamil \emph{maruvu}; -\te{కో} is the reflexive ending.
\end{cousinweb}

\begin{grammarlens}[title={What you've built}]
One more word for describing what you lived and felt.
\end{grammarlens}

\subsection*{Your turn}
\begin{itemize}
\raggedright
  \item \emph{Say it:} \emph{meccukō}
  \item \emph{Say it:} \emph{meccukō}, once more
  \item \emph{Say it:} \emph{nānna nā bommanu meccukunnāru.} — Dad praised my picture.
  \item \emph{Recall:} say the Telugu for to admire, then the Telugu for an admirer, then say \emph{meccukō} again
\end{itemize}

\begin{tcolorbox}[breakable,colback=teal!4,colframe=teal!35!black,title={Before you move on}]
What does \te{మెచ్చుకో} mean? (\textquotedblleft{}To appreciate; to praise\textquotedblright{}.) Say it once more.
\end{tcolorbox}
\section[\texorpdfstring{\te{ప్రశంస} (praśaṁsa)}{ప్రశంస (praśaṁsa)}]{\te{ప్రశంస} (praśaṁsa) — acclaim; praise}

\label{lesson:TE-C504-prasamsa}



\begin{quote}
Before the new one: say the Telugu for an admirer, then the Telugu for to appreciate.
\end{quote}

\subsection*{You'll want to know: \te{ప్రశంస}}
\textbf{\te{ప్రశంస}} — \emph{praśaṁsa} — \textquotedblleft{}acclaim; praise\textquotedblright{}.

\begin{quote}
\te{ఆ} \te{సినిమాకు} \te{చాలా} \te{ప్రశంసలు} \te{వచ్చాయి}.
\end{quote}

\begin{quote}
\emph{ā sinimāku cālā praśaṁsalu vaccāyi.}
\end{quote}

\begin{quote}
That film won a lot of acclaim.
\end{quote}

Plural: \te{ప్రశంసలు}, \emph{praśaṁsalu}.

Often plural: \te{ప్రశంసలు} \te{అందుకో}, to receive praise.

\begin{cousinweb}
\textbf{\te{ప్రశంస}} (\emph{praśaṁsa}) is borrowed from Sanskrit \emph{praśaṁsā}, \textquotedblleft{}praise, commendation, fame\textquotedblright{} (Monier-Williams).
\end{cousinweb}

\begin{grammarlens}[title={What you've built}]
One more word for describing what you lived and felt.
\end{grammarlens}

\subsection*{Your turn}
\begin{itemize}
\raggedright
  \item \emph{Say it:} \emph{praśaṁsa}
  \item \emph{Say it:} \emph{praśaṁsa}, once more
  \item \emph{Say it:} \emph{ā sinimāku cālā praśaṁsalu vaccāyi.} — That film won a lot of acclaim.
  \item \emph{Recall:} say the Telugu for an admirer, then the Telugu for to appreciate, then say \emph{praśaṁsa} again
\end{itemize}

\begin{tcolorbox}[breakable,colback=teal!4,colframe=teal!35!black,title={Before you move on}]
What does \te{ప్రశంస} mean? (\textquotedblleft{}Acclaim; praise\textquotedblright{}.) Say it once more.
\end{tcolorbox}
